\subsection{کالبدشکافی ماژول \lr{simple\_ssm}}

\subsubsection*{۱. فلسفه طراحی، پارادایم معماری و تمثیل ملموس}
ماژول مینیمال \lr{simple\_ssm} جایگزینی سبک و کاملاً خطی برای مکانیزم توجه سنتی است که فرآیند مدل‌سازی زبانی را در قالب معادلات دیفرانسیل بازگشتی خطی گسسته‌شده پیاده می‌کند و حافظه استنتاج را در حد یک بردار ثابت نگه می‌دارد.

\textbf{تمثیل ملموس و مثال کاربردی:}
مدار فیلتر پایین‌گذر پسیو \lr{RC} در مهندسی الکترونیک: ولتاژ ورودی پیوسته بار الکتریکی خازن را تغییر می‌دهد و ولتاژ خروجی به صورت تابعی نرم و زوال‌یابنده از تاریخچه بارهای الکتریکی گذشته ظاهر می‌شود، بدون آنکه مدار نیازی به ذخیره‌سازی داده‌های تک‌تک پالس‌های گذشته در یک حافظه دیجیتال داشته باشد.

\subsubsection*{۲. فرمول‌بندی و استنتاج عمیق ریاضی \lr{(Mathematical Derivation)}}
مدل فضای حالت با معادلات بازگشتی گسسته زیر فرمول‌بندی می‌شود:
\begin{equation}
\mathbf{h}_t = \bar{\mathbf{A}} \mathbf{h}_{t-1} + \bar{\mathbf{B}} \mathbf{x}_t, \quad \mathbf{y}_t = \mathbf{C} \mathbf{h}_t + \mathbf{D} \mathbf{x}_t
\end{equation}
که در آن ماتریس انتقال حالت $\mathbf{A}$ قطری بوده و با تبدیل دوخطی \lr{(Bilinear Transform)} به فرم گسسته درمی‌آید:
\begin{equation}
\bar{\mathbf{A}} = \exp(\boldsymbol{\Delta}_t \mathbf{A}), \quad \bar{\mathbf{B}} = (\boldsymbol{\Delta}_t \mathbf{A})^{-1} (\bar{\mathbf{A}} - \mathbf{I}) \cdot (\boldsymbol{\Delta}_t \mathbf{B})
\end{equation}
محاسبه کل دنباله در فاز آموزش با اسکن موازی پیشوندی \lr{(Parallel Scan)} در زمان $\mathcal{O}(N)$ اجرا می‌گردد.

\subsubsection*{۳. کالبدشکافی الگوریتمی و جریان داده \lr{(Algorithmic Tensor Trace)}}
در پیاده‌سازی سورس پایتون در فایل ماژول \lr{simple\_ssm.py}، پردازش در متد \lr{forward} به صورت مرحله‌به‌مرحله و گام‌به‌گام صورت می‌پذیرد:
\begin{enumerate}
\item دریافت بردار ورودی $x_t$ و محاسبه گام زمانی پیوسته $\Delta_t$.
\item گسسته‌سازی ماتریس‌های قطری فضای حالت با عملگر نمایی پایدار عددی.
\item اجرای اسکن پیشوندی موازی بر روی دنباله زمانی جهت انباشت حالت‌های مخفی.
\item پروجکشن خطی حالت‌ها با ماتریس خروجی $\mathbf{C}$ و ترکیب با مسیر مستقیم $\mathbf{D}$.
\item نرمال‌سازی خروجی لایه با عملگر پیوسته \lr{LayerNorm} و ارسال به لایه بعد.
\end{enumerate}

\subsubsection*{۴. تحلیل همجوشی سخت‌افزاری و ترافیک حافظه \lr{(Hardware \& Roofline Analysis)}}
از دیدگاه مدل سقف کارایی \lr{(Roofline Model)} و معماری پردازنده‌های گرافیکی مدرن (نظیر کارت گرافیک \lr{NVIDIA RTX 4070}):
\begin{itemize}
\item \textbf{ثابت بودن مطلق بافر استنتاج $\mathcal{O}(1)$:} نیازی به بافر ذخیره کلیدها و مقادیر \lr{(KV-Cache)} وجود ندارد و مصرف رم در طول زمان همواره تخت می‌ماند.
\item \textbf{حداکثر کارایی حافظه رجیسترها:} اسکن خطی در حافظه کش اشتراکی \lr{SRAM} محصور شده و به \lr{DRAM} مراجعه نمی‌کند.
\end{itemize}

\subsubsection*{۵. ارزیابی تجربی، بنچمارک و توجیه علمی نتایج}
بر اساس آزمون‌های ثبت‌شده در خط لوله ارزیابی پایان‌نامه:
\begin{itemize}
\item \textbf{وضعیت رسمی آزمون کیفی:} \lr{PASS}
\item \textbf{دقت ثبت‌شده روی بنچ‌مارک:} \lr{62.13\%}
\item \textbf{شاخص ارزیابی آماری } \lr{(F1 Score)}\rl{:}  \lr{63.8\%}
\item \textbf{اوج مصرف حافظه گرافیکی } \lr{(Peak VRAM)}\rl{:}  \lr{620.0 MB} (در برابر \lr{۴۲۰۰ MB} مدل مبنا)
\item \textbf{نرخ توان عملیاتی } \lr{(Throughput)}\rl{:}  \lr{44,500} توکن بر ثانیه
\end{itemize}

\textbf{توجیه علمی و علت‌شناسی جذاب نتایج:}
دستیابی به بالاترین نرخ توان عملیاتی در این رده یعنی \lr{44,500} توکن در ثانیه و مصرف فوق‌العاده اندک \lr{620.0 MB} حافظه (کاهش \lr{85.2\%}ی نسبت به \lr{۴۲۰۰ MB} مبنا) با دقت \lr{62.13\%}، نشان می‌دهد که \lr{SimpleSSM} گزینه‌ای بی‌رقیب برای سناریوهای با منابع سخت‌افزاری به شدت محدود و پردازش فوری جریان داده‌های مداوم است.

\subsubsection*{۶. دیاگرام معماری شکیل در لاتک \lr{(TikZ Architecture Flowchart)}}
جهت تبیین دقیق جریان داده تانسوری و نحوه تغییر شکل ماتریس‌ها، نمودار معماری این ماژول با استفاده از بسته قدرتمند \lr{TikZ} در محیط لاتک ترسیم شده است:

\begin{center}
\begin{tikzpicture}[scale=0.82, every node/.style={scale=0.82}, 
    node distance=1.1cm,
    auto,
    block/.style={rectangle, draw=blue!70, fill=blue!12, rounded corners=4pt, text centered, minimum height=0.85cm, font=\small},
    ssm/.style={rectangle, draw=orange!80, fill=orange!18, rounded corners=4pt, text centered, minimum height=0.85cm, font=\small},
    outblock/.style={rectangle, draw=teal!80, fill=teal!15, rounded corners=4pt, text centered, minimum height=0.85cm, font=\small}
]
    \node [block] (input) {ورودی $x_t$};
    \node [ssm, right=1.0cm of input] (state) {حالت پنهان $h_t$};
    \node [ssm, right=1.0cm of state] (scan) {اسکن خطی قطری};
    \node [outblock, right=1.0cm of scan] (out) {خروجی خطی};

    \draw [->, thick, draw=blue!80] (input) -- (state);
    \draw [->, thick, draw=orange!80] (state) -- (scan);
    \draw [->, thick, draw=orange!80] (scan) -- (out);
\end{tikzpicture}
\end{center}

\subsubsection*{۷. پیوست تصویری و پرامپت تولید هوش مصنوعی}
برای ترسیم تجسم سه‌بعدی و گرافیک صنعتی این مکانیزم با ابزارهای هوش مصنوعی، پرامپت انگلیسی زیر با دقت تدوین شده است:

\vspace{0.2cm}
\begin{latin}
\footnotesize
\textbf{AI Technical Visualization Prompt:}\\
"A 3D minimalist scientific diagram of SimpleSSM. A linear harmonic oscillator model where data pulses smoothly modulate an internal glowing resonant core with zero memory overhead, clean clinical white and neon blue lighting, 8k resolution --ar 16:9"
\end{latin}

\begin{figure}[H]
    \centering
    \framebox[0.9\textwidth]{
    \begin{minipage}{0.85\textwidth}
    \vspace{1.8cm}
    \centering
    \textbf{طرح‌واره معماری و جریان تانسوری ماژول \lr{simple\_ssm}}\\[0.3cm]\small \lr{(Tensor Flow \& Architectural Schematic for simple\_ssm)}\\[0.2cm]
    \small (بر اساس پرامپت فنی فوق تولید و در این کادر جایگزین می‌شود)
    \vspace{1.8cm}
    \end{minipage}
    }
    \caption{شماتیک و جریان داده تانسوری در ماژول \lr{simple\_ssm}}
    \label{fig:simple_ssm}
\end{figure}

\vspace{0.5cm}
\hrule
\vspace{0.5cm}
